% !TeX program = xelatex
% Copyright (C) 2026 bsxiaocai
% SPDX-License-Identifier: LPPL-1.3c
% Example distributed under LPPL 1.3c. See LICENSE and NOTICE.md.
% Compile from the project root: xelatex -output-directory=build examples/cool-demo.tex
% 逐页操作说明：docs/templates/cool.md；全部尺寸：docs/customization.md。
\documentclass[aspectratio=169]{ctexbeamer}
\usetheme{HUAS}
\HUASsetup{style=cool,logo=true,footer=standard,sectionpage=true,
  faculty={教学与学术交流}}
\input{examples/local-assets.tex}
% Cool 页眉和页脚的总高度 = height + depth；下面均为可取消注释的示例。
% \HUASsetlayout{headline width=\paperwidth,headline height=3mm,
%   headline depth=1.3mm,footline height=3.5mm,footline depth=1.3mm,
%   cool title padding=0.8em,frametitle padding=2mm}
% 封面主标题/副标题/作者/机构/日期分别在以下五行替换。
\title{阅读、证据与教学设计}
\subtitle{HUAS-Beamer 冷色样式与学术内容组件}
\author{HUAS-Beamer 项目}
\institute{湖南文理学院}
\date{\today}
\begin{document}
\HUASmaketitle
\begin{frame}{内容结构}
  \tableofcontents[hideallsubsections]
\end{frame}
\section{研究与论证}
\begin{frame}{结论、证据与边界}
  \begin{block}{核心判断}
    阅读任务应让学生给出文本证据，而不只是复述教师结论。
  \end{block}
  \begin{alertblock}{适用边界}
    单个教学案例不能直接证明所有学生都适用同一策略。
  \end{alertblock}
  \begin{exampleblock}{可以观察的表现}
    学生能定位语句、解释推理，并回应不同理解。
  \end{exampleblock}
\end{frame}
\begin{frame}{引文与出处}
  \begin{HUASquotation}{《论语·为政》}
    学而不思则罔，思而不学则殆。
  \end{HUASquotation}
  \medskip
  引文与讨论分开呈现，读者可以辨认原文和汇报者的分析。
\end{frame}
\begin{frame}{研究流程：问题如何变为证据}
  \begin{enumerate}
    \item \textbf{提出问题}：明确需要解释的阅读现象。
    \item \textbf{收集材料}：保留文本批注与课堂讨论记录。
    \item \textbf{比较解释}：寻找支持与反对判断的证据。
    \item \textbf{说明边界}：交代样本范围与尚未解决的问题。
  \end{enumerate}
\end{frame}
\section{文本与教学}
\begin{frame}{文本细读：原文与分析并列}
  % 左右各 48% 正文宽，默认各 6.72cm；引文框随左栏 linewidth 缩放。
  \begin{columns}[T,onlytextwidth]
    \begin{column}{0.48\textwidth}
      \begin{HUASquotation}{陆游《游山西村》}
        山重水复疑无路，\\柳暗花明又一村。
      \end{HUASquotation}
    \end{column}
    \begin{column}{0.48\textwidth}
      \begin{block}{细读问题}
        \begin{itemize}
          \item 哪些词语表现视野的变化？
          \item 两句之间如何形成转折？
          \item 解释如何回到具体语词？
        \end{itemize}
      \end{block}
    \end{column}
  \end{columns}
\end{frame}
\begin{frame}{教学案例：目标、任务与反馈}
  \begin{block}{学习目标}
    学生能从一段文本中选择证据，解释自己的理解。
  \end{block}
  \begin{exampleblock}{任务与反馈}
    \begin{itemize}
      \item 先独立标注，再与同伴比较不同的解释。
      \item 汇报时展示语句与推理，反馈聚焦证据是否充分。
    \end{itemize}
  \end{exampleblock}
\end{frame}
\begin{frame}{表格：把任务和证据对应起来}
  \begin{table}
    \caption{教学活动与可观察证据}
    % 1.4 是表格行高倍数；lll 为三列左对齐，改列宽时检查总宽度。
    \renewcommand{\arraystretch}{1.4}
    \begin{tabular}{lll}
      \hline
      \textbf{环节} & \textbf{阅读任务} & \textbf{证据形式} \\
      \hline
      初读 & 记录疑问 & 问题清单 \\
      细读 & 标注关键语句 & 带说明的批注 \\
      讨论 & 比较不同解释 & 观点与依据 \\
      反思 & 修改最初判断 & 修订记录 \\
      \hline
    \end{tabular}
  \end{table}
  \medskip
  表格适合比较关系；长段原文应拆到文本细读页。
\end{frame}
\begin{frame}{图片：交代对象与来源}
  \begin{figure}
    \IfFileExists{assets/research/HUAS_building_2_bgcolor.png}{%
      % 正文图宽是正文宽的 65%（默认 9.1cm），不受建筑装饰区高度参数影响。
      \includegraphics[width=0.65\textwidth]{assets/research/HUAS_building_2_bgcolor.png}%
      \caption{校园建筑线稿}%
    }{%
      % 缺图占位框高 2.2cm、内部宽为正文宽的 60%；caption 在两分支各自设置。
      \fbox{\parbox[c][2.2cm][c]{0.6\textwidth}{\centering 校园图片}}%
      \caption{图片位置示例}%
    }
  \end{figure}
  {\usebeamerfont{HUAS source}\color{HUASMuted}
    \IfFileExists{assets/research/HUAS_building_2_bgcolor.png}{%
      图源：项目所有者提供的本地研究素材。%
    }{图片缺失，当前显示占位框。}\par}
  \medskip
  图注说明对象，来源说明帮助读者追溯素材。
\end{frame}
\begin{frame}[fragile]{代码：用短片段说明处理过程}
  % verbatim 必须保留 fragile；small 仅缩小当前 block 内的代码字号。
  \begin{block}{Python 示例：统计已经完成的阅读任务}
    \small
    \begin{verbatim}
tasks = ["read", "annotate", "discuss"]
completed = {"read", "annotate"}
progress = sum(task in completed for task in tasks)
print(progress, "/", len(tasks))
    \end{verbatim}
  \end{block}
  \medskip
  代码页使用原生 verbatim；中文解释放在代码区外。
\end{frame}
\section{总结与资料}
% keypoint 自己生成普通内容页并计入页码，第二参数支持手动换行。
\HUASkeypoint{关键结论}{让问题引导阅读，\\让证据支持判断。}
\begin{frame}{参考文献与进一步阅读}
  \begin{thebibliography}{9}
    % 手工条目；{9} 是编号宽度样本，超过 9 条改成 {99}。
    \bibitem{classic}《论语·为政》.
      \newblock 本示例引文的出处。
    \bibitem{beamer}Till Tantau, Joseph Wright, Vedran Miletić.
      \newblock The beamer class: User Guide.
      \newblock TeX Live 随附的 Beamer 手册。
    \bibitem{huas}HUAS-Beamer 项目.
      \newblock 使用说明与阶段开发记录.
      \newblock 本仓库 README 与 docs。
  \end{thebibliography}
\end{frame}
\HUASclosingframe
\end{document}
